\chapter{The geometric inequality I: unitary towers, hermitian lattices and slopes}\label{ch:geoa}

This chapter covers the objects behind the archive's ``geometric inequality'' and the
form in which that inequality reaches the analytic part of the proof. It describes the
ranks and hermitian spaces of a tower of unitary groups, the abstract theta chain, the
hermitian lattices over a number field with their slopes, a slope inequality of Bost
type and a frame theorem built on it, and the abstract Arakelov interface \texttt{DV.Geo}.
It then states the slope record, a list of real numbers with inequalities between them,
together with the window hypotheses and the tie that says a record comes from a family
of packages. Downstream, the analytic chain uses only the real-number clauses of a tied
record and three arithmetic identifications; the archimedean and finite sides and the Chebyshev inequality come next.

\section{Ranks and the tower of hermitian spaces}



\begin{lemma}[rank identities]\label{geoa:rank-identities}
\lean{Formalism.BasicNotions.Ranks.rS_two, Formalism.BasicNotions.rank_base_space, Formalism.BasicNotions.factXm, Formalism.BasicNotions.factYm}
\uses{host:ranks}
For all ranks: $2r_S+2=m$; $N(0)+1+2r_S=r$; $N(m)=\operatorname{rk}X'+1$; and
$\operatorname{rk}Y(m)=\operatorname{rk}X'$.
\end{lemma}
\begin{proof}
\uses{host:ranks}
Integer arithmetic from the definitions, using that $m$ is even and $m\ge2$. The last
identity holds by definition.
\end{proof}

\begin{definition}[hermitian spaces with base lattice]\label{geoa:hspace}
\lean{Formalism.BasicNotions.CMData, Formalism.BasicNotions.HSpace}
A \emph{CM datum} for number fields $F\subset E$ (with an involution on $E$) records that
$E$ is a CM field unramified in the sense of the archive's \texttt{IsCMUnram}, that $F$ is
totally real, that $[E:F]=2$, and that the fixed field of the involution is $F$.
A \emph{hermitian space} over $E$ is a finite-dimensional $E$-vector space $V$ with a
nondegenerate hermitian form $h$, a natural number $\rk$ equal to $\dim_EV$, a basis,
a function assigning to each embedding $\tau\colon E\to\CC$ a pair of natural numbers
(the declared signature), and an $\mathcal O_E$-submodule $L_0\subset V$ that is a
lattice and is unimodular for $h$. It satisfies \texttt{SigOK} if, for every $\tau$, the
Gram matrix in the given basis has signature $\operatorname{sig}(\tau)$ after applying
$\tau$; it is \emph{definite} if for every $\tau$ the signature is $(\rk,0)$ or
$(0,\rk)$. An $E$-linear map $\iota$ is an \emph{isometric embedding} if
$h'(\iota x,\iota y)=h(x,y)$ for all $x,y$.
\end{definition}

\begin{definition}[the spaces of a choice of ranks]\label{geoa:spaces}
\lean{Formalism.BasicNotions.Spaces}
\uses{host:ranks, geoa:hspace}
Let $R$ be ranks and $F\subset E$ number fields. \emph{Spaces} for $R$ consist of a CM
datum for $E/F$, a prime number $T$, and hermitian spaces over $E$, all satisfying
\texttt{SigOK}, as follows.
\begin{itemize}
\item $X_j$ for $j\in\N$, with $\rk X_j=N(j)$ and $X_j$ definite for $j\le m$.
\item $X'$ of rank $\operatorname{rk}X'$, an isometric embedding $X'\to X_m$, and a
vector $z\in X_m$ with $h(z,z)\ne0$ orthogonal to the image of $X'$.
\item $Y_j$ for $j\in\N$, with $\rk Y_j=\operatorname{rk}Y(j)$ for $j\le m$; linear maps
$X_j\to Y_j$ which are isometric embeddings for $j<m$; an $E$-linear isomorphism
$X'\cong Y_m$ that is isometric.
\item In $Y_0$: a vector $z_m$ with $h(z_m,z_m)=-T$, orthogonal to the image of $X_0$, and
vectors $e_1,\dots,e_{r_S},f_1,\dots,f_{r_S}$ with $h(e_i,e_{i'})=h(f_i,f_{i'})=0$ and
$h(e_i,f_{i'})=\delta_{ii'}$, orthogonal to the image of $X_0$ and to $z_m$; and
$\operatorname{sig}Y_0(\tau)=\operatorname{sig}Y_0$ for every $\tau$.
\item $W_j$ for $j\in\N$, with $\rk W_j=r$ and signature $\operatorname{sig}W(j)$ at
every $\tau$ for $1\le j\le m$; for each such $j$ there are a bijection of the basis
index sets and $g\in\GL_r(E\otimes_F\A_F^{\mathrm f})$ with $g^*G_{Y_0}g=G_{W_j}$, where
$G$ denotes the Gram matrix and $\A_F^{\mathrm f}$ the finite adeles.
\item $V_j$ for $j\in\N$, with $\rk V_j=n_V$ for $1\le j\le m$, a sign
$\epsilon_j\in\{\pm1\}$, and $L_0(V_j)$ an orthonormal lattice of rank $n_V$ for
$\epsilon_jh$.
\end{itemize}
\end{definition}

\begin{lemma}[definiteness at the top]\label{geoa:top-definite}
\lean{Formalism.BasicNotions.factXmDef}
\uses{geoa:spaces, geoa:rank-identities}
For spaces of ranks $R$, the spaces $X'$ and $Y_m$ are definite, and for every $\tau$
the space $X_m$ has signature $(N(m),0)$ at $\tau$ if and only if $X'$ has signature
$(\operatorname{rk}X',0)$ at $\tau$.
\end{lemma}
\begin{proof}
\uses{geoa:spaces, geoa:rank-identities}
The isometric embedding $X'\to X_m$ identifies the Gram matrix of $X'$ with the
restriction of the form of $X_m$, which is definite. A restriction of a definite form
to a nonzero subspace is definite of the same sign; the rank identity
$N(m)=\operatorname{rk}X'+1$ and the CM datum (embeddings commute with the
involution) give the signature statement. The isometry $X'\cong Y_m$ transfers
definiteness to $Y_m$.
\end{proof}

\begin{remark}[the top of the tower]\label{geoa:rem-top}
$\rk Y_j-\rk X_j=(r+jt)-N(j)=m-1-j$, so an isometric embedding of the nondegenerate
space $X_j$ into $Y_j$ can exist only for $j\le m-1$, which is where the archive asks for
one. At $j=m$ the direction reverses: $Y_m\cong X'$ embeds into $X_m$, whose rank is one
larger, and $z$ spans the complementary line. The archive states that $z$ is anisotropic
and orthogonal to the image, but not that $X_m$ is the orthogonal sum. Likewise, the
vectors $z_m,e_i,f_i$ in $Y_0$ are only required to be orthogonal to $X_0$ and to each
other; that they span a complement of $X_0$ follows from the rank identity
$N(0)+1+2r_S=r$, which the definition does not state.
\end{remark}

\begin{definition}[column chain and theta kernels]\label{geoa:column-chain}
\lean{Formalism.BasicNotions.ColumnChain, Formalism.BasicNotions.ColumnChain.Kmap, Formalism.BasicNotions.ColumnChain.K, Formalism.BasicNotions.ColumnChain.col, Formalism.ThetaLiftIndependence.kerOpM, Formalism.ThetaLiftIndependence.chainOpM}
\uses{host:ranks}
Let $R$ be ranks, $\iota$ a finite set and $(\mathrm{hc}_{\tau})_{\tau\in\iota}$
Harish-Chandra parameters for $R$. A \emph{column chain} consists of, for each
$i\in\N$: a finite set $\operatorname{cls}_i$; for each $c\in\operatorname{cls}_i$ a
finite group $\operatorname{stab}_i(c)$; a type $Y_i$; nowhere-zero functions
$\operatorname{de}_i\colon Y_i\times\iota\to\CC$; functions $\rho_i$ on
$\operatorname{cls}_i$ and $\rho'_i$ on $\operatorname{cls}_{i+1}$ of absolute value $1$
for $i<m$; $\CC$-linear maps
$\operatorname{unf}_i(c,c')\colon\CC^{Y_i}\to\CC^{Y_{i+1}}$; integers
$s^{\sharp}_i(\tau),s^{\mathrm{div}}_i(\tau)$ given for $i<m$ by explicit affine
expressions in the parameter $e$ of $\mathrm{hc}_{\tau}$; and a positive integer
$D_{\Gamma}$ divisible by every $\#\operatorname{stab}_i(c)$, $i\le m$. Put
$\operatorname{wt}_i(c)=1/\#\operatorname{stab}_i(c)$ and
\[
K_i(c,c')=\rho_i(c)\rho'_i(c')\cdot M_{i+1}^{-1}\circ\operatorname{unf}_i(c,c')\circ
M^{\sharp}_i,
\]
where $M^{\sharp}_i$ multiplies by $y\mapsto\prod_{\tau}\operatorname{de}_i(y,\tau)^{s^{\sharp}_i(\tau)}$
and $M_{i+1}$ by $y'\mapsto\prod_{\tau}\operatorname{de}_{i+1}(y',\tau)^{s^{\mathrm{div}}_i(\tau)}$.
A \emph{form at level $i$} is a function $\varphi\colon\operatorname{cls}_i\to\CC^{Y_i}$.
The kernel operator is $(K_i\varphi)(c')=\sum_{c}\operatorname{wt}_i(c)K_i(c,c')\varphi(c)$,
and the \emph{column} is the $m$-fold composite $\operatorname{col}=K_{m-1}\circ\cdots\circ
K_0$.
\end{definition}

\begin{remark}[the column chain is abstract]\label{geoa:rem-column-abstract}
Nothing in the column chain refers to the spaces of the previous definitions. The sets
$\operatorname{cls}_i$, the groups $\operatorname{stab}_i(c)$, the types $Y_i$ and the maps
$\operatorname{unf}_i$ are free fields. Reading $\operatorname{cls}_i$ as the class set of
the genus of $L_0(X_i)$, the weights as mass weights, and $K_i$ as theta kernels for
$U(X_i)\times U(X_{i+1})$ is an interpretation that this structure does not enforce;
any such link must come from other fields of a package.
\end{remark}

\begin{definition}[columns at $T$]\label{geoa:columns-at-T}
\lean{Formalism.BasicNotions.ColumnsAtT, Formalism.BasicNotions.ColumnsAtT.Mram, Formalism.BasicNotions.ColumnsAtT.du}
\uses{geoa:column-chain}
For a column chain, \emph{columns at $T$} consist of a prime $T$; a finite set
$v_T$ with $f$ elements; for each $v\in v_T$ a finite field $F_v$ with $T$ elements, a
nontrivial additive character $\psi_v$ of $F_v$ and a homomorphism $\rho_T(v)$ from
$F_v^{\times}$ to the endomorphisms of the space of forms at level $0$; a
finite-dimensional subspace $M_{u_0}$ of forms at level $0$ stable under every
$\rho_T(v)(c)$; and constants $d^*_u\in\N_{>0}$, $D^*>0$, $\Gamma_P>1$. Put
$\operatorname{Nm}_v=\sum_{c\in F_v^{\times}}\rho_T(v)(c)$,
$M_{\mathrm{ram}}=M_{u_0}\cap\bigcap_v\ker\operatorname{Nm}_v$ and
$d_u=\dim_{\CC} M_{\mathrm{ram}}$.
\end{definition}

\begin{definition}[period map]\label{geoa:period-map}
\lean{Formalism.BasicNotions.PeriodMap, Formalism.BasicNotions.PeriodMap.Ptil, Formalism.BasicNotions.PeriodMap.PtilCol}
\uses{geoa:column-chain}
A \emph{period map} for a column chain consists of complex vector spaces $M_t$, $W_z$,
$W_{\mathrm{row}}$, linear maps $\operatorname{ev}_z$ from forms at level $m$ to $W_z$ and
$\operatorname{low}\colon W_z\to W_{\mathrm{row}}$, and a bilinear pairing
$\operatorname{per}\colon W_{\mathrm{row}}\times M_t\to\CC$. For a form $\varphi$ at level
$0$ and $F\in M_t$ put
$\widetilde P(\varphi,F)=\operatorname{per}(\operatorname{low}(\operatorname{ev}_z(\operatorname{col}\varphi)),F)$.
All spaces and maps are free fields.
\end{definition}

\section{Hermitian lattices and slopes}

\begin{definition}[hermitian lattices]\label{geoa:hlat}
\lean{Formalism.BasicNotions.HLat, Formalism.BasicNotions.HLat.Good, Formalism.BasicNotions.HLat.rank, Formalism.HermitianLatticesTheory.fcoord, Formalism.HermitianLatticesTheory.realize}
Let $K$ be a number field. A \emph{hermitian lattice} over $K$ is a $K$-vector space $V$
with a basis $b$ indexed by a finite set $\iota$, a complex matrix $X_v\in M_{\iota}(\CC)$
for each infinite place $v$ of $K$, and an $\mathcal O_K$-submodule $L\subset V$. It is
\emph{good} if every $\det X_v$ is nonzero, $L$ is finitely generated and $L$ spans $V$.
Its rank is $\dim_K(K L)$. The norm at $v$ is
$\|x\|_v=|X_v^{-1}\sigma_v(x_b)|$, the euclidean norm of $X_v^{-1}$ applied to the
coordinate vector of $x$ under the embedding $\sigma_v$ attached to $v$. The
\emph{realization} sends $x$ to the real and imaginary parts of all these vectors, a
group homomorphism from $V$ to a euclidean space of dimension $2\cdot\#\{v\}\cdot\#\iota$.
\end{definition}

\begin{definition}[heights and slopes]\label{geoa:slopes}
\lean{Formalism.HN.ht, Formalism.HN.muDef, Formalism.HN.muMaxDef, Formalism.HN.muMinDef, Formalism.HN.opNormAt, Formalism.HN.htMap, Formalism.HN.slope, Formalism.HN.muMax, Formalism.HN.muMin, Formalism.F.covol, Formalism.F.slope, Formalism.BasicNotions.HLat.mu, Formalism.BasicNotions.HLat.muMax, Formalism.BasicNotions.HLat.muMin}
\uses{geoa:hlat}
For a $\Z$-submodule $\Lambda$ of a euclidean space, $\operatorname{covol}\Lambda$ is the
square root of the Gram determinant of a $\Z$-basis if $\Lambda$ is discrete and $0$
otherwise; $\deg\Lambda=-\log\operatorname{covol}\Lambda$ and
$\operatorname{sl}(\Lambda)=\deg\Lambda/\rk_{\Z}\Lambda$. For an $\mathcal O_K$-submodule $L$
of a hermitian lattice write $\Lambda(L)$ for its image under the realization, $n_{\infty}$
for the number of infinite places of $K$, and $c_K=-\log2+\log|d_K|/[K:\Q]$. Then
\begin{align*}
\operatorname{ht}(L)&=2^{n_{\infty}\rk L}\,|d_K|^{-\rk L/2}\operatorname{covol}\Lambda(L),\\
\mu(L)&=2\operatorname{sl}\Lambda(L)+c_K,\\
\mu_{\max}(L)&=2\sup\{\operatorname{sl}\Lambda(N):0\ne N\le L\}+c_K,\\
\mu_{\min}(L)&=2\inf\{\operatorname{sl}(\Lambda(L)/\Lambda(N)):N\le L\text{ saturated},\ N\ne L\}+c_K.
\end{align*}
For a $K$-linear map $\varphi\colon V\to W$ between hermitian lattices, with matrix
$\Phi$ in the given bases, put
$\|\varphi\|_v=\|X'^{-1}_v\,\sigma_v(\Phi)\,X_v\|$ (operator norm) and
\[
\operatorname{ht}(\varphi)=\frac1{n_{\infty}}\sum_v\log\|\varphi\|_v^2 .
\]
\end{definition}

\begin{remark}[normalisations]\label{geoa:rem-normalisation}
The matrices $X_v$ are frames: the metric at $v$ is the one making the columns of $X_v$
orthonormal. Nothing in the definition of a hermitian lattice requires a hermitian form
or positivity beyond invertibility. The slope is twice the euclidean slope of the
realization plus the constant $c_K$, and the height of a map uses only archimedean
operator norms (squared, averaged over the infinite places). There is no finite-place
contribution; integrality enters separately, as an inclusion of lattices. The suprema
and infima are of real sets, so they take Mathlib's junk value when the set is empty or
unbounded.
\end{remark}

\begin{definition}[duals]\label{geoa:dual}
\lean{Formalism.HN.dualLat, Formalism.HN.dualFrame, Formalism.BasicNotions.HLat.dualL, Formalism.BasicNotions.HLat.dualB, Formalism.BasicNotions.HLat.dualX, Formalism.BasicNotions.HLat.muDual, Formalism.BasicNotions.HLat.muMaxDual, Formalism.BasicNotions.HLat.muMinDual}
\uses{geoa:hlat, geoa:slopes}
The dual of a hermitian lattice $A$ is the space $V^*$ with the dual basis, the frames
$(X_v^{T})^{-1}$, and the lattice $L^{\vee}=\{\lambda\in V^*:\lambda(L)\subset\mathcal O_K\}$.
Its slope, maximal slope and minimal slope are written $\mu(A^{\vee})$,
$\mu_{\max}(A^{\vee})$ and $\mu_{\min}(A^{\vee})$.
\end{definition}

\begin{definition}[period lattice data]\label{geoa:divdata}
\lean{Formalism.BasicNotions.DivData, Formalism.BasicNotions.sFrame, Formalism.BasicNotions.uFrame, Formalism.BasicNotions.DivData.cc, Formalism.BasicNotions.DivData.OneMapIntoCC, Formalism.BasicNotions.DivData.HBound}
\uses{geoa:hlat, geoa:slopes, geoa:dual}
Let $K$ be a number field and $\mathcal W$ a finite set (the words). \emph{Period
lattice data} consist of a hermitian lattice $R$ (test forms); for each $u\in\mathcal W$
a hermitian lattice $N_u$ and a $K$-linear map $B_u\colon V(N_u)\to V(R)^*$; reals
$h_u$, $HH$, $\eta_F$, $\delta_P$ with $\delta_P\le1$; a natural number $D_{\mathrm{all}}$;
and the inequality $s+\mu_{\max}(R)\le HH$, where
\[
s=\frac{\sum_u d_u(\mu(N_u)-h_u)}{\sum_ud_u},\qquad d_u=\rk N_u .
\]
If $\mathcal W$ is nonempty, $u_F=\max_u(\mu_{\max}(N_u)-h_u)-s$. Put $cc=R^{\vee}$. The
data satisfy \texttt{OneMapIntoCC} if $B_u(N_u)\subset cc$ for all $u$, and
\texttt{HBound} if $\operatorname{ht}(B_u)\le h_u$ for all $u$, the height taken from
the frames of $N_u$ to the dual frames of $R$.
\end{definition}

\begin{remark}\label{geoa:rem-divdata}
Integrality and the height bound are separate predicates, not fields: period lattice
data alone do not force $B_u$ into $R^{\vee}$, and $B_u$ need not be injective. The
inequality $s+\mu_{\max}(R)\le HH$ is a field. A package (\texttt{D7}) carries one set of
period lattice data together with invertibility of all the frames
(\texttt{div\_good}). When $\sum_ud_u=0$ Lean's division gives $s=0$.
\end{remark}

\begin{theorem}[slope inequality]\label{geoa:slope-ineq-mean}
\lean{Formalism.HN.SlopeIneqMean, Formalism.HN.slope_inequality_mean}
\uses{geoa:slopes}
Let $K$ be a totally complex number field, $V,W$ finite-dimensional $K$-vector spaces with
bases and invertible frames at every infinite place, $\varphi\colon V\to W$ a $K$-linear
map and $A\subset V$ a finitely generated $\mathcal O_K$-submodule with $\varphi(A)\ne0$.
Then:
\begin{enumerate}
\item $\mu_{\min}(A)-\operatorname{ht}(\varphi)\le\mu_{\min}(\varphi(A))$;
\item for every finitely generated $B\subset W$ containing $\varphi(A)$,
$\mu_{\min}(A)\le\mu_{\max}(B)+\operatorname{ht}(\varphi)$, and if moreover $\varphi$ is
injective on $A$, then $\mu_{\max}(A)\le\mu_{\max}(B)+\operatorname{ht}(\varphi)$.
\end{enumerate}
\end{theorem}
\begin{proof}
% uses-lean: Formalism.HN.slopeIneqMean_of, Formalism.HN.twist_of, Formalism.HN.rescale_covolume_formula, Formalism.HN.place_lattice_structure, Formalism.HN.slope_inequality_general
The archive derives it from a general slope inequality for euclidean lattices and
group homomorphisms (\texttt{slope\_inequality\_general}), transported to
$\mathcal O_K$-lattices by the structure of the realization (rank of $\Lambda(L)$ is
$[K:\Q]\rk L$), the rescaling formula for covolumes under scalar changes of frames, and
shift invariance of the slopes under such twists. The proof is in
\arch{Formalism/HarderNarasimhan/HarderNarasimhanProofs/} and is not expanded here.
\end{proof}

\begin{remark}[relation with Bost's inequality]\label{geoa:rem-bost}
Since $\mu(A)\le\mu_{\max}(A)$, the last clause gives
$\mu(A)\le\mu_{\max}(B)+\operatorname{ht}(\varphi)$ for $\varphi$ injective on $A$ with
$\varphi(A)\subset B$. This is Bost's slope inequality in the form with integral maps and
archimedean norms only, in the archive's normalisation of $\mu$. The archive does not
apply it directly to the period maps $B_u$; it is used inside the proof of the frame
theorem below.
\end{remark}

\begin{lemma}[frame linear algebra]\label{geoa:frame-la1}
\lean{Formalism.HN.FrameLA1, Formalism.HN.frame_linear_algebra_first, Formalism.CorankShareLayerCake.FR.hs2, Formalism.CorankShareLayerCake.FR.frK, Formalism.CorankShareLayerCake.FR.frQ}
Let $(V_i)_{i\in I}$ be a finite family of finite-dimensional complex inner product spaces,
$H$ another, $B_i\colon V_i\to H$ linear, $c_i>0$, $N_i\subset V_i$ subspaces, and
$e,\eta,\varepsilon$ real. Write $\|A\|_{\mathrm{HS}}^2=\operatorname{tr}(A^*A)$ and
$Q=\sum_ic_iB_iB_i^*$. Assume $\sum_i\dim V_i>0$,
$\sum_i\|c_iB_i^*B_i-1\|_{\mathrm{HS}}^2\le e\sum_i\dim V_i$, $0\le e\le1/4$, $\eta\ge0$,
$\operatorname{tr}(Q^2)\le(1+\eta)(\operatorname{tr}Q)^2/\dim H$ and
$\sum_i(\dim V_i-\dim N_i)\le\varepsilon\sum_i\dim V_i$. Then
\[
\dim\Bigl(\sum_iB_i(N_i)\Bigr)^{\perp}\le(\eta+4\varepsilon+4\sqrt e)\dim H .
\]
\end{lemma}
\begin{proof}
Finite-dimensional linear algebra in
\arch{Formalism/HarderNarasimhan/HarderNarasimhanProofs/FrameLinearAlgebraFirst.lean};
not expanded here.
\end{proof}

\begin{theorem}[frame theorem, sharp form]\label{geoa:theoremF-sharp}
\lean{Formalism.HN.TheoremFSharp, Formalism.HN.frame_theorem_sharp, Formalism.HN.theoremFSharp_of}
\uses{geoa:slope-ineq-mean, geoa:frame-la1, geoa:slopes}
Let $K$ be a totally complex number field. Let $W$ have a basis indexed by a finite set
$\kappa$, invertible frames $X'_w$, and a finitely generated $\mathcal O_K$-submodule
$cc$ spanning $W$. Let $I$ be a nonempty finite set and, for $i\in I$, let $V_i$ have a
basis indexed by $\iota_i$, invertible frames $X_{i,w}$, a finitely generated
$\mathcal O_K$-submodule $N_i$ spanning $V_i$, a $K$-linear map $B_i\colon V_i\to W$ and
a real $h_i$. Fix an infinite place $v$, a finite set $\kappa_1$, invertible matrices
$Y_i\in M_{\iota_i}(\CC)$, $Y'\in M_{\kappa}(\CC)$, matrices
$\mathrm{Ball}_i\in M_{\kappa\sqcup\kappa_1,\iota_i}(\CC)$, weights $\gamma_i$, and reals
$s,u,HH,e,\eta,\delta$. Assume:
\begin{enumerate}
\item $\sum_i\#\iota_i>0$; $B_i(N_i)\subset cc$; and $\operatorname{ht}(B_i)\le h_i$
whenever $B_i\ne0$;
\item $s\sum_i\#\iota_i\le\sum_i\#\iota_i\,(\mu(N_i)-h_i)$; $u\ge0$ and
$\mu_{\max}(N_i)-h_i\le s+u$ for all $i$; $s-\mu_{\min}(cc)\le HH$ and $HH\ge0$;
\item $0\le\delta<1$ and $(1-\delta)\#(\kappa\sqcup\kappa_1)\le\#\kappa$;
\item the $\kappa$-rows of $\mathrm{Ball}_i$ are those of
$Y'^{-1}\sigma_v(\text{matrix of }B_i)Y_i$;
\item $\gamma_i>0$, $\sum_i\|\gamma_i\mathrm{Ball}_i^*\mathrm{Ball}_i-1\|_{\mathrm{HS}}^2\le
e\sum_i\#\iota_i$, $0\le e\le1/4$, and, with
$Q=\sum_i\gamma_i\mathrm{Ball}_i\mathrm{Ball}_i^*$,
$\operatorname{tr}(Q^2)\le(1+\eta)(\operatorname{tr}Q)^2/\#(\kappa\sqcup\kappa_1)$ and
$\eta\ge0$.
\end{enumerate}
Then
\[
s-(1-\delta)^{-1}\Bigl(4u\log\Bigl(1+\frac{HH}u\Bigr)+(\eta+4\sqrt e)HH\Bigr)\le\mu(cc).
\]
\end{theorem}
\begin{proof}
\uses{geoa:slope-ineq-mean, geoa:frame-la1}
% uses-lean: Formalism.HN.thmF_share, Formalism.HN.thmF_place, Formalism.HN.hnStep
For $\tau\ge0$ let $N_i(\tau)\subset N_i$ be the Harder--Narasimhan step of $N_i$ made of
the part of slope at least $s-\tau+h_i$ (in the normalisation of $\mu$), and $S(\tau)$
the step of $cc$ of slope at least $s-\tau$. A counting lemma on the slope lists
(\texttt{thmF\_share}, using the mean bound $s$ and the cap $s+u$) gives
$\sum_i(\#\iota_i-\rk N_i(\tau))\le\frac u{u+\tau}\sum_i\#\iota_i$. The first clause of
the slope inequality, applied to $B_i$ on $N_i(\tau)$ with the height bound, shows that
$B_i(N_i(\tau))$ has minimal slope at least $s-\tau$, hence lies in $S(\tau)$. The frame
linear algebra at the place $v$, with $\varepsilon=u/(u+\tau)$ and hypothesis~(3), gives
$\#\kappa-\rk\sum_iB_i(N_i(\tau))\le w(\tau)\#\kappa$ with
$w(\tau)=(1-\delta)^{-1}(\eta+4u/(u+\tau)+4\sqrt e)$, so the same bound holds for
$\#\kappa-\rk S(\tau)$. Expressing $\mu(cc)$ through its Harder--Narasimhan polygon,
using $\mu_{\min}(cc)\ge s-HH$, and integrating the rank bound over $\tau\in[0,HH]$ gives
the stated inequality, since $\int_0^{HH}4u/(u+\tau)\,d\tau=4u\log(1+HH/u)$.
\end{proof}

\begin{theorem}[frame theorem]\label{geoa:theoremF}
\lean{Formalism.HN.TheoremF, Formalism.HN.frame_theorem_log_ideal, Formalism.HN.theoremF_of_sharp, Formalism.HN.theoremF}
\uses{geoa:theoremF-sharp}
Under the hypotheses of the sharp form,
\[
s-\log\#I-(1-\delta)^{-1}\Bigl(4u\log\Bigl(1+\frac{HH}u\Bigr)+(\eta+4\sqrt e)HH\Bigr)\le\mu(cc).
\]
\end{theorem}
\begin{proof}
\uses{geoa:theoremF-sharp}
Since $I$ is nonempty, $\log\#I\ge0$; subtract it from the left side of the sharp form.
\end{proof}

\begin{remark}[on the frame theorem]\label{geoa:rem-theoremF}
The term $-\log\#I$ is a weakening of the sharp form that the archive carries along; the
slope clause of the record below has the same term with $\#I$ the number of words. The
rank weights are the basis sizes $\#\iota_i$, which equal $\rk N_i$ because $N_i$ spans
$V_i$. The matrices $Y_i$, $Y'$ are arbitrary invertible matrices, not tied to the frames
$X$; the proof uses them only through ranks. For $u=0$ Lean reads $HH/u$ as $0$, and the
term $4u\log(1+HH/u)$ is $0$. The survey describes this theorem as conditioned on
\texttt{FrameLA1}; in the archive \texttt{FrameLA1} is proved
(\texttt{frame\_linear\_algebra\_first}).
\end{remark}

\section{The abstract Arakelov interface}

\begin{definition}[geometry interface]\label{geoa:geo}
\lean{Formalism.BasicNotions.DV.Geo}
\uses{geoa:hlat}
Let $F$ be a number field. A \emph{geometry} over $F$ is a structure with: a type
$\mathrm{Var}$ of ``varieties'' $Y$ with dimensions; function fields $\operatorname{Fn}(Y)$
(fields with an $F$-algebra structure); abelian groups $\operatorname{Dv}(Y)$ of
divisors with an effectivity predicate closed under $0$ and sums; spaces
$H^0(Y,\Delta)\subset\operatorname{Fn}(Y)$, finite-dimensional, multiplicative in $\Delta$
and monotone under effective differences; a canonical divisor; closed-subset types with
properness, inclusion and support; an intersection number of $\dim Y$ divisors; predicates
$\operatorname{nef}$, $\operatorname{big}$, base-point freeness, birationality and
non-contraction; vanishing predicates; complex points with a measure for each embedding
$F\to\CC$; types of metrics with sum, power and Fubini--Study constructions, pointwise
norms and positivity predicates; a type of adelic line bundles with a construction
$\operatorname{ofW}$ from a hermitian lattice and a map to $\operatorname{Fn}(Y)$,
twists, an arithmetic top degree, an arithmetic volume and a nefness predicate; sup norms
at infinite and finite places with the expected behaviour under $\operatorname{ofW}$ and
twists; and several cut and residue predicates. The only axioms are the effectivity,
$H^0$ and sup-norm compatibilities just listed.
\end{definition}

\begin{definition}[the one-point geometry]\label{geoa:pointgeo}
\lean{GlobalTheory.PointGeometry.pointGeo, GlobalTheory.PointGeometry.AdelicLine, GlobalTheory.PointGeometry.AdelicLine.deg}
\uses{geoa:geo}
For a number field $F$, the one-point geometry has $\mathrm{Var}=\{\ast\}$ of dimension
$0$, $\operatorname{Fn}(\ast)=F$, one divisor, every divisor effective,
$H^0=F$ for every divisor, intersection number $1$, and
\[
\operatorname{nef}:=\text{True},\qquad \operatorname{big}:=\text{True}.
\]
Its complex points are a single point with the Dirac measure, metrics are nonnegative
reals, and its adelic line bundles are \emph{adelic lines}: positive reals at infinite
places and at almost all finite places equal to $1$, with arithmetic degree
$-\bigl(\sum_vm_v\log\ell_v+\sum_w\log\ell_w\bigr)/[F:\Q]$; the arithmetic volume is the
positive part of the degree.
\end{definition}

\begin{lemma}[a geometry with a nonzero section exists]\label{geoa:exists-geo}
\lean{GlobalTheory.PointGeometry.exists_geo, GlobalTheory.PointGeometry.point_global_sections_ne_bot}
\uses{geoa:pointgeo}
For every number field $F$ there are a geometry $G$ over $F$ and $Y\in\mathrm{Var}(G)$ with
$H^0(Y,0)\ne0$.
\end{lemma}
\begin{proof}
\uses{geoa:pointgeo}
Take the one-point geometry; $H^0(\ast,0)=F\ne0$.
\end{proof}

\begin{definition}[divisor data of a package]\label{geoa:divx}
\lean{Formalism.BasicNotions.DivX, Formalism.BasicNotions.DV.hExpl6, Formalism.BasicNotions.DV.KapRBase}
\uses{geoa:geo, geoa:hspace, geoa:hlat}
% uses-lean: Formalism.BasicNotions.D7, Formalism.BasicNotions.SizeData₂
For a package $D$ with coefficient field $LL$, \emph{divisor data} consist of a geometry
over $LL$; a level type and a family of ``Shimura data'' attached to hermitian spaces over
the host field (abstract records of Chern and arithmetic degree functionals); for each
prime $p$ a nonempty finite set of places above $p$ with nonnegative weights summing to
$1$ given by ramification indices and norms; integral models at $p$ with first Chern
class data; for each level $j$ a good hermitian lattice $M_j$ embedded isometrically in the
package's space of forms on $U(W_j)$, and maps from $M_j$ into
$H^0(Y,\Lambda_j)$ for varieties $Y$ of the geometry; local twist data; positive reals
$\kappa_R(e)$ for embeddings $e\colon LL\to\CC$, invariant under conjugation; and
surface invariants. \texttt{KapRBase} says that $\kappa_R$ at the base embedding equals
the explicit constant $\kappa^{\sharp}(D)$. The explicit height $h^{(6)}_u$ of a word $u$ is
an explicit combination of size data of $D$ minus the mean of $\log\kappa_R(e)$ over the
embeddings; it uses no other field of the divisor data.
\end{definition}

\begin{lemma}[provisional divisor data]\label{geoa:provisional-divx}
\lean{GlobalTheory.exists_provisionalDivX, GlobalTheory.provisionalDivX}
\uses{geoa:divx, geoa:exists-geo}
Let $D$ be a package such that for each $1\le j\le m$ the space of forms $M_j$ has a
complex basis in which every generator of the given lattice has coordinates in the
image of $LL$. Then there are divisor data $X$ for $D$ satisfying \texttt{KapRBase}.
\end{lemma}
\begin{proof}
\uses{geoa:exists-geo}
% uses-lean: GlobalTheory.WittFracComplex.nonempty_ringHom_complex, GlobalTheory.nonempty_admissibleModel
Take the geometry and variety supplied by the existence lemma, which is the one-point
geometry, together with an embedding of the Witt-vector fraction field into $\CC$,
admissible integral models, and lattices $M_j$ built from the given bases; set
$\kappa_R\equiv\kappa^{\sharp}(D)$ and take the surface types empty.
\end{proof}

\begin{remark}[the geometry interface is hollow on this path]\label{geoa:rem-hollow-geo}
\texttt{DV.Geo} is not built on Mathlib schemes: it is a record of types and predicates
whose only axioms are bookkeeping compatibilities. It is instantiated by the one-point
model with $\operatorname{nef}:=\text{True}$ and $\operatorname{big}:=\text{True}$, and
the archive's provisional divisor data use exactly this model. Two further instances in
\arch{Formalism/LocalToGlobal/LatticeGeometry/Inhabited.lean}
(\texttt{geoDegenerate}, \texttt{geoDegenerateNoTau}) also set $\operatorname{nef}$ and
$\operatorname{big}$ to True, and \texttt{geoExample} has an empty type of varieties.
A family stores divisor data and defines the heights $h_u$ by $h^{(6)}_u$, which reads
only $\kappa_R$. The lemma \texttt{exists\_provisionalDivX} itself is not reached from the final
theorems, but the one-point model is: a review traced, link by link, the chain
\texttt{hypotheses\_nonempty\_of\_unramified\_exponents\_alt},
\texttt{hyps\_e\_third\_layer\_ser\_n\_of\_geometric\_galois},
\texttt{exists\_hyp\_n\_of\_geometric\_galois\_compat}, \dots,
\texttt{exists\_divisor\_eigen\_index\_third\_component\_alt},
\texttt{ThirdIdx.nonempty\_divX}
(\arch{Development/AbelianVarietyTheory/Proofs/IwasawaDecomposition/ExistsEigenReadThirdLocal.lean}, lines 24--26), whose last
step builds \texttt{provisionalDivX} from \texttt{exists\_geo}. So divisor data over
the one-point model are used on the main chain. Whether later steps use any geometric
content of those data was not traced. Statements about nef or big divisors made through this interface carry no
geometric content unless the instance is a real one.
\end{remark}

\section{The inequalities on real numbers}

\begin{definition}[slope records]\label{geoa:record}
\lean{Formalism.RankGapBounds.WeightedMeanData.Data, Formalism.RankGapBounds.WeightedMeanData.Consts}
A \emph{slope record} with index type $U_0$ (shifts), word types $W(u_0)$ for
$u_0\in U_0$, a type $\mathrm{Top}$ (top places) and natural numbers $m,f$ consists of
real numbers $x,\eps,a,b,b',k,L,L_a,B_0,N_c,\gamma,\ell_K,v,\lg,\log\Gamma_L$; reals
$a_{\tau}$ for $\tau\in\mathrm{Top}$; weights $w(u_0,j,\rho)$ for $j<m$, $\rho<f$; for
$u\in W(u_0)$ reals $d(u),T(u),\mu_N(u),\mu_N^{\max}(u),h(u),\mu_M(u),\mu_M^{\max}(u)$
and $E_j(u)$ for $j<m$; for $u_0\in U_0$ reals $\mu_R(u_0),\mu_R^{\min}(u_0)$,
$\mathrm{Fo},\mathrm{Ga},\mathrm{Ta},\mathrm{PS},\mathrm{IO},HH,\delta,\beta_1,\beta_2$
(Lean fields \texttt{muCC}, \texttt{muMinCC}, \texttt{Fo}, \texttt{Ga}, \texttt{Ta},
\texttt{ASK}, \texttt{IO}, \texttt{HH}, \texttt{del}, \texttt{$\eta$}, \texttt{e}); and reals
$\bar\delta,\bar\beta_1,\bar\beta_2$. \emph{Constants} are ten reals
$c_{\mathrm{Div}},e_{\mathrm{Div}},c_E,c_H,c_{\mathrm{Count}},c_{\mathrm{Slope}},
c_{\mathrm{Sl}},c_{\mathrm{Loss}},e_{\mathrm{Loss}},c_{HH}$.
\end{definition}

\begin{definition}[derived quantities of a record]\label{geoa:record-derived}
\lean{Formalism.RankGapBounds.WeightedMeanData.ph, Formalism.RankGapBounds.WeightedMeanData.wmean, Formalism.RankGapBounds.WeightedMeanData.umean, Formalism.RankGapBounds.WeightedMeanData.Data.nT, Formalism.RankGapBounds.WeightedMeanData.Data.sumA, Formalism.RankGapBounds.WeightedMeanData.Data.logOm, Formalism.RankGapBounds.WeightedMeanData.Data.Ee, Formalism.RankGapBounds.WeightedMeanData.Data.s, Formalism.RankGapBounds.WeightedMeanData.Data.uF, Formalism.RankGapBounds.WeightedMeanData.Data.LoP, Formalism.RankGapBounds.WeightedMeanData.Data.Main}
\uses{geoa:record}
For a slope record with $U_0$ and all $W(u_0)$ finite and nonempty, put
$n_T=\#\mathrm{Top}$, $\Sigma a=\sum_{\tau} a_{\tau}$, $\operatorname{ph}(y)=y\log y-y$,
$\log\Omega=\gamma+\log(\pi/v)-\ell_K$, $\operatorname{Ee}(u)=\sum_{j<m}E_j(u)$,
\[
s(u_0)=\frac{\sum_ud(u)(\mu_N(u)-h(u))}{\sum_ud(u)},\qquad
u_F(u_0)=\max_u(\mu_N^{\max}(u)-h(u))-s(u_0),
\]
$\mathrm{Lo}_P=c_{\mathrm{Loss}}n_Tmb\,x^{e_{\mathrm{Loss}}\eps}$ and
\[
\operatorname{Main}(u_0)=\sum_{j<m}\sum_{\tau} a_{\tau}\Bigl(\log a_{\tau}-1-\gamma+\ell_K
-\frac1f\sum_{\rho<f}\log w(u_0,j,\rho)\Bigr).
\]
\end{definition}

\begin{definition}[clauses of the size data, first group]\label{geoa:clauses}
\lean{Formalism.RankGapBounds.WeightedMeanData.Data.HidaSetting, Formalism.RankGapBounds.WeightedMeanData.Data.DIV, Formalism.RankGapBounds.WeightedMeanData.Data.HComp, Formalism.RankGapBounds.WeightedMeanData.Data.Slopes, Formalism.RankGapBounds.WeightedMeanData.Data.SlopeIneq, Formalism.RankGapBounds.WeightedMeanData.Data.HHBound, Formalism.RankGapBounds.WeightedMeanData.Data.FParams}
\uses{geoa:record-derived}
For a record and constants, quantified over all $u_0$ and $u,u'\in W(u_0)$:
\begin{itemize}
\item (positivity, \texttt{HidaSetting}) $x,b,k,v>0$, $f>0$, $a,b'\ge0$, $a_{\tau}\ge0$,
$w>0$, $d>0$, $T>0$, $\mu_N\le\mu_N^{\max}$, $c_{\mathrm{Sl}}\ge0$, $\bar\delta<1$,
$\bar\beta_1\ge0$, $0\le\bar\beta_2\le1/4$;
\item (\texttt{DIV}) $\mu_R+h(u)-\mathrm{Fo}\le\sum_{\tau}\bigl(m(\operatorname{ph}(a_{\tau})
-a_{\tau}\log\Omega)-a_{\tau}\log T(u)\bigr)+c_{\mathrm{Div}}n_TaL_a
+c_{\mathrm{Div}}m(B_0+N_c)x^{e_{\mathrm{Div}}\eps}$;
\item (\texttt{HComp}) $\bigl|(h(u)+\Sigma_{\tau} a_{\tau}\log T(u))-(h(u')+\Sigma_{\tau}
a_{\tau}\log T(u'))\bigr|\le c_H(n_TaL_a+mB_0)$;
\item (\texttt{Slopes}) $-c_{\mathrm{Slope}}n_TaL_a(1+m/k)\le\mu_M(u)$ and
$\mu_M^{\max}(u)\le c_{\mathrm{Slope}}n_TaL_a$;
\item (\texttt{SlopeIneq}) if $0\le\delta<1$, $\beta_1\ge0$, $0\le\beta_2\le1/4$,
$s-\mu_R^{\min}\le HH$ and $HH>0$, then
\[
s-\log\#W(u_0)-(1-\delta)^{-1}\Bigl(4u_F\log\Bigl(1+\frac{HH}{u_F}\Bigr)
+(\beta_1+4\sqrt{\beta_2})HH\Bigr)\le\mu_R;
\]
\item (\texttt{HHBound}) $s-\mu_R^{\min}\le HH$, $HH>0$ and $HH\le x^{c_{HH}}b$;
\item (\texttt{FParams}) $0\le\delta\le\bar\delta$, $0\le\beta_1\le\bar\beta_1$,
$0\le\beta_2\le\bar\beta_2$.
\end{itemize}
\end{definition}

\begin{definition}[clauses of the size data, second group]\label{geoa:clauses-5}
\lean{Formalism.LFunctionChainParameters.WeightedSlopeBounds.Lo5, Formalism.LFunctionChainParameters.WeightedSlopeBounds.QuotientSlopeBoundsWith, Formalism.LFunctionChainParameters.WeightedSlopeBounds.EComp5, Formalism.LFunctionChainParameters.WeightedSlopeBounds.Count5, Formalism.LFunctionChainParameters.WeightedSlopeBounds.LossBounds}
\uses{geoa:record-derived}
Let $X''$ be real and put
$\mathrm{Lo}_5(u)=c_{\mathrm{Sl}}\mathrm{IO}+c_{\mathrm{Sl}}n_Tm(a/k+b')\log T(u)
+c_{\mathrm{Sl}}X''+1$. For all $u_0$ and $u,u'\in W(u_0)$:
\begin{itemize}
\item (\texttt{EComp5}) $|\operatorname{Ee}(u)-m\,\Sigma a\,\lg|\le c_EmB_0$;
\item (\texttt{Count5}) $\sum_{j<m}f^{-1}\sum_{\rho}\log(w(u_0,j,\rho)v/\pi)+m\lg
-c_{\mathrm{Count}}(L_a+m^2L_a/k)\le\log T(u)\le c_{\mathrm{Count}}mL$ and
$|\log T(u)-\log T(u')|\le c_{\mathrm{Count}}\log k$;
\item (\texttt{QuotientSlopeBoundsWith}) $\mu_N^{\max}(u)\le\mu_M^{\max}(u)+\mathrm{Ga}
+\mathrm{Fo}-\operatorname{Ee}(u)+c_{\mathrm{Sl}}mB_0$ and
$\mu_M(u)+\mathrm{Ga}+\mathrm{Fo}-\operatorname{Ee}(u)-\mathrm{Lo}_5(u)\le\mu_N(u)$;
\item (\texttt{LossBounds}) $\mathrm{PS}(u_0)-\mathrm{Ga}(u_0)\le\mathrm{Lo}_P$ and
$\mathrm{IO}(u_0)\le\mathrm{Lo}_P$.
\end{itemize}
\end{definition}

\begin{definition}[size data of a record]\label{geoa:sizedata}
\lean{Formalism.LFunctionChainParameters.LChain.SizeData₂}
\uses{geoa:clauses, geoa:clauses-5}
A record \emph{satisfies the size data} with constants $C$ and real $X''$ if it satisfies
\texttt{HidaSetting}, \texttt{DIV}, \texttt{EComp5}, \texttt{HComp}, \texttt{Count5},
\texttt{Slopes}, \texttt{QuotientSlopeBoundsWith} (with $X''$), \texttt{LossBounds},
\texttt{SlopeIneq}, \texttt{HHBound} and \texttt{FParams}.
\end{definition}

\begin{remark}[names and unused variants]\label{geoa:rem-variants}
The name \texttt{HidaSetting} denotes a list of positivity conditions; an unrelated
predicate \texttt{DV.HidaSetting} on hermitian spaces exists. The name \texttt{SizeData}$_2$
is also used for the structure \texttt{Formalism.BasicNotions.SizeData}$_2$ of a package,
which is a different object. The file of the record also defines \texttt{EComp}
(a bound for each $j$), \texttt{Count} (with $L$ in place of $L_a$ and an extra term in
$\log\Gamma_L$), \texttt{SlopesN}, \texttt{Losses} (bounds for the mean over $u_0$) and
the proposition \texttt{MeanBoundSecondForm}; none is a clause of the size data, and no
proof of \texttt{MeanBoundSecondForm} was found. The field $\log\Gamma_L$ enters no clause
of the size data. The hypotheses of \texttt{SlopeIneq} follow from \texttt{HidaSetting},
\texttt{HHBound} and \texttt{FParams}.
\end{remark}

\begin{remark}[the slope clause and the frame theorem]\label{geoa:rem-slopeineq}
\texttt{SlopeIneq} has the shape of the conclusion of the frame theorem with $I=W(u_0)$,
target $cc=R^{\vee}$, sources $N_u$, $s$ the rank-weighted mean of $\mu(N_u)-h_u$,
$u=u_F$, $\eta=\beta_1$ and $e=\beta_2$. It bounds the mean $s$ by the slope of
$R^{\vee}$, not by its maximal slope, which is stronger than Bost's inequality applied word
by word. Together with \texttt{DIV}, an upper bound for $\mu_R+h(u)-\mathrm{Fo}$, it
pins $\mu_R$ from both sides. In the record the clause is a hypothesis on real numbers.
\end{remark}

\begin{definition}[window hypotheses]\label{geoa:errwindows}
\lean{Formalism.LFunctionChainParameters.LErr.ConstsBound, Formalism.LFunctionChainParameters.LErr.ErrWindows0, Formalism.LFunctionChainParameters.LErr.ErrWindows}
\uses{geoa:record}
Let $X'',e_K,K_W$ be real. A record with constants satisfies \texttt{ErrWindows} if:
$1\le x$; $0<\eps\le1$; $1\le b$; $xb\le a$; $1\le n_T$; $m\le x^{1/2+\eps/2}$;
$xm\le k$; $0\le e_K$; $1\le k\le x^{e_K}$; $k\le a$; $0\le L\le2x^{\eps}$;
$0\le L_a\le2x^{\eps}$; $0\le b'$ and $kb'\le a$; $0\le X''\le a$; $0\le a_{\tau}\le2a$;
$B_0,N_c\ge0$; $\log\#W(u_0)\le(k-3)\log2$ for all $u_0$; $\bar\delta\le1/2$;
$\bar\beta_1,\bar\beta_2\ge0$; the constants $c_{\mathrm{Div}},c_E,c_H,c_{\mathrm{Count}},
c_{\mathrm{Slope}},c_{\mathrm{Sl}},c_{\mathrm{Loss}}$ lie in $[0,K_W]$ and
$e_{\mathrm{Div}},e_{\mathrm{Loss}},c_{HH}\ge0$; and the defect bound
$(\bar\beta_1+4\sqrt{\bar\beta_2})\,x^{c_{HH}}b\le K_W\,x^{\eps}\,n_Ta$.
\end{definition}

\begin{definition}[identification and size bounds]\label{geoa:ident}
\lean{Formalism.LFunctionChainParameters.LChain.Ident, Formalism.LFunctionChainParameters.LChain.BNBounds, Formalism.LFunctionChainParameters.ThickSetMeanBounds.ValuesAtX}
\uses{geoa:record}
% uses-lean: Formalism.LFunctionChainParameters.characterSumData.ArithData, Formalism.LFunctionChainParameters.characterSumData.ASK, Formalism.LFunctionChainParameters.characterSumData.GPsi, Formalism.LFunctionChainParameters.characterSumData.chC, Formalism.D5.PhiSet
\uses{bal:additive-model, bal:arith-data, bal:ask, bal:gpsi} % resolved
Let $H$ be a finite abelian group, $\mathcal D$ arithmetic data on $H$, $N\in\N$,
$\ell$ a function on the characters of $E_0\times H$, $TT\subset H$, $a\colon H\to\R$,
$w\colon H\times\{0,\dots,m-1\}\times(E_0\times H)\to\R$, and bijections
$\mathrm{Top}\cong TT$ and $\{0,\dots,f-1\}\cong\Phi$. A record with $U_0=H$ satisfies
\texttt{Ident} if its $a_{\tau}$ and $w$ are carried to $a$ and $w$ by these bijections,
$\gamma$ is Euler's constant, $\ell_K=\Re\ell(\chi_C)$, and
$\mathrm{PS}(u_0)=\sum_{t\in TT}\sum_{j<m}\operatorname{ASK}(a_t,t,w(u_0,j))$ for all
$u_0$. For $T_2\subset H$ it satisfies \texttt{BNBounds} if
$B_0\le2\#H\,m\,b\,x^{\eps}$ and $N_c\le\#T_2\,k\,b\,x^{\eps}/(\#H\cdot\#H\,x^{-\eps})$.
For real $x,\eps,C_L$, \texttt{ValuesAtX} holds if
$|\Psi_{\theta}(x)|\le x^{1/2+\eps}$ and $|\ell(\theta)|\le C_L\log x$ for the three
characters $\theta\in\{\psi_C,\chi_C,\chi_C\psi_C\}$, where $\Psi_{\theta}$ is the
archive's twisted prime-sum function \texttt{GPsi}.
\end{definition}

\begin{lemma}[the main term]\label{geoa:main-eq}
\lean{Deduction.ch_top_main_eq}
\uses{geoa:record-derived, geoa:ident}
% uses-lean: Formalism.LFunctionChainParameters.characterSumData.MAIN, Deduction.CH.IdentU
\uses{bal:main-defect} % resolved
Let a record, $\mathcal D,N,\ell,TT,a,w$ and the bijections be as above, let
$\operatorname{tr}\colon U_0\to H$, and assume the variant \texttt{IdentU} of
\texttt{Ident} in which $w(u_0,j,\cdot)$ is matched with $w(\operatorname{tr}u_0,j,\cdot)$
(\texttt{Ident} is the case $U_0=H$, $\operatorname{tr}=\mathrm{id}$). Then for every
$u_0$,
\[
\operatorname{Main}(u_0)=\sum_{t\in TT}\sum_{j<m}\Re\,\operatorname{MAIN}(a_t,w(\operatorname{tr}u_0,j)),
\]
where $\operatorname{MAIN}(a,w)=a\bigl(\log a-1-\gamma+\ell(\chi_C)-\#\Phi^{-1}\sum_{\rho\in\Phi}\log w(\rho)\bigr)$.
\end{lemma}
\begin{proof}
\uses{geoa:record-derived}
The bijection $\{0,\dots,f-1\}\cong\Phi$ gives $f=\#\Phi$ and turns sums over $\rho<f$
into sums over $\Phi$. Substituting the identifications of $a_{\tau}$, $w$, $\gamma$ and
$\ell_K$ into $\operatorname{Main}$, exchanging the sums over $j$ and $\tau$ and
reindexing $\tau$ by $TT$ gives the right side term by term, since all quantities other
than $\ell(\chi_C)$ are real. No positivity of $a_t$ or $w$ is used.
\end{proof}

\begin{definition}[families]\label{geoa:famcore}
\lean{Formalism.MainNotions.FamCore, Formalism.MainNotions.FamP}
\uses{geoa:divx}
% uses-lean: Formalism.BasicNotions.D7, Formalism.BasicNotions.SizeData₂
Let $\eps$ be real. A \emph{core family} at $\eps$ consists of a threshold $x_1>1$; an
index type $\mathrm{Ix}(x)$ for each real $x$; an unbounded set $\mathrm{Win}\subset\R$
with $\mathrm{Ix}(x)$ nonempty for $x\in\mathrm{Win}$ and $\mathrm{Ix}(0)$ nonempty;
packages $D(x,i)$ with package size data $L(x,i)$ at scale $x$ (for $x>x_1$) and $\eps$;
uniformity across indices of ranks, $TT$, $T_2$, $q$, $d$, $Q_S$, Artin numerics, the
auxiliary primes and further constants; bounds by a constant $C_o$; a common word type;
and a covering clause realising every residue $u_0<n$ by some index. A \emph{family}
extends a core family by parameter choices, tie data, divisor data $X(x,i)$, geometric
instance data, explicit definitions of $\nu$, $N'$ and $R$, and the height formula
$h_u=h^{(6)}_u(D(x,i),L(x,i),X(x,i),C_{\mathrm{sz}})$, among other fields.
\end{definition}

\begin{definition}[uniform explicit-formula condition]\label{geoa:efamc}
\lean{Formalism.HC.EAt, Formalism.HC.EFamC}
\uses{geoa:famcore}
% uses-lean: Formalism.HC.EAtH
For a level $\Lambda$ and a core family, \texttt{EFamC} holds if there are $C$ and $x_0$
such that for all $x\ge x_0$ and all indices $i$, the host of $D(x,i)$ satisfies the
explicit-formula condition $\mathrm{EAtH}(\Lambda,TT,x,\eps,C)$.
\end{definition}

\begin{definition}[tied chains]\label{geoa:tied}
\lean{Formalism.MainNotions.TiePts, Formalism.MainNotions.TieGlobal, Formalism.MainNotions.TieLocal, Formalism.MainNotions.TieArith, Formalism.MainNotions.ChainTied}
\uses{geoa:famcore, geoa:sizedata, geoa:errwindows, geoa:ident, geoa:divdata, geoa:dual}
% uses-lean: Formalism.LFunctionChainParameters.LChain.GrowthInequalityRemainingHyps, Formalism.LFunctionChainParameters.LChain.CostWindows, Formalism.LFunctionChainParameters.LChain.RestSize, Formalism.MainNotions.arithOf
\uses{bal:rest-size} % resolved
Let $\mathrm{Fm}$ be a family at $\eps$ and $e_D,e_L,K_W,c_{HH},e_K,C_L,x$ real.
$\mathrm{Fm}$ is \emph{tied} at $x$ if there exist an index $i_0$, $m,f\in\N$, a choice
$\operatorname{pt}(u_0)$ of index for each $u_0$ in the group $A$ of the frame of
$D(x,i_0)$, integer sequences $a^v,b^v$, a finite type $\mathrm{Top}$, a record with
$U_0=A$ and $W(u_0)$ the words of $D(x,\operatorname{pt}(u_0))$, constants, a real $X''$,
and $\ell,N,TT,T_2,a,w$, bijections and reals $\eta,X^{\flat},\theta$, such that:
\begin{enumerate}
\item (\texttt{TiePts}) all $D(x,\operatorname{pt}(u_0))$ have the frame of $D(x,i_0)$,
the given $m$, $f$, $q$, $d$, shift $u_0$, and $a_t,b_t$ read off from $a^v,b^v$;
\item (\texttt{TieGlobal}) $a,b,b',k,L,L_a,B_0,N_c$ are those of the size data, $\gamma$ is
Euler's constant, $\ell_K$ is the recorded logarithmic derivative at $1$, $v$ is the
host's $v$, $\lg=\log(g_0g_1)$, $\log\Gamma_L$, $X''$, $\eta$ and $X^{\flat}$ are as
recorded;
\item (\texttt{TieLocal}) $d(u)=\rk N_u$, $T(u)$ is the level prime of the word,
$\mu_N(u),\mu_N^{\max}(u)$ are the slope and maximal slope of $N_u$, $h(u)=h_u$,
$\mu_M,\mu_M^{\max}$ are those of the form lattice of $u$, $E_j(u)$, $\mathrm{Ta}$ and
$\mathrm{IO}$ are sums of local sizes over the support of the size data,
$\mu_R=\mu(R^{\vee})$, $\mu_R^{\min}=\mu_{\min}(R^{\vee})$, $\mathrm{Fo},\mathrm{Ga}$ are the
size-data fields, and $\delta$ is the package's rank defect;
\item (\texttt{TieArith}) $TT$, $T_2$, $a$ and $w$ are those of the members, and
$\ell(\psi)=L'(1,\chi)/L(1,\chi)$ for every nontrivial $\psi$ and every primitive
Dirichlet character $\chi$ agreeing with $\psi$ on positive integers;
\item $x,\eps$ are those of the record, $e_{\mathrm{Div}}\le e_D$, $e_{\mathrm{Loss}}\le e_L$,
$c_{HH}^{\mathrm{rec}}\le c_{HH}$;
\item the record satisfies \texttt{Ident}, the size data, \texttt{ErrWindows},
\texttt{BNBounds}, the growth-inequality hypotheses (including \texttt{ValuesAtX}),
\texttt{CostWindows} and \texttt{RestSize}.
\end{enumerate}
The fields $HH$, $\beta_1$, $\beta_2$ and $\mathrm{PS}$ are set by no tie;
$\mathrm{PS}$ is fixed by \texttt{Ident}.
\end{definition}

\begin{lemma}[ties are discarded downstream]\label{geoa:tied-core}
% overlaps (deduplicated): geob:tiedcore
\lean{Formalism.MainNotions.ChainTied.core_of}
\uses{geoa:tied}
% uses-lean: Formalism.LFunctionChainParameters.LChain.ChainCore, UpperTheory.Comparison₂.FamCoreConstantBounds, Formalism.LFunctionChainParameters.characterSumData.HostFacts
\uses{bal:split-outside, geob:coreath} % resolved
If $\mathrm{Fm}$ is tied at $x$ and every host of the family satisfies \texttt{HostFacts},
then \texttt{ChainCore} holds at $x$: there exist a finite abelian group, a record over it
and the remaining data satisfying \texttt{Ident}, the size data, \texttt{ErrWindows},
\texttt{BNBounds}, the growth-inequality hypotheses, \texttt{CostWindows} and
\texttt{RestSize}. Also, if $\mathrm{Fm}$ is tied at $x$ then the core family satisfies
\texttt{FamCoreConstantBounds} at $x$.
\end{lemma}
\begin{proof}
\uses{geoa:tied}
Unpack the existential. The first statement drops \texttt{TiePts}, \texttt{TieGlobal},
\texttt{TieLocal} and \texttt{TieArith} and repackages the rest with the group $A$. The
second also drops \texttt{TiePts}, \texttt{TieGlobal} and \texttt{TieLocal}, keeps from
\texttt{TieArith} only the descriptions of $TT$ and $T_2$ and the values of $\ell$, and
identifies $TT$ with the host's set.
\end{proof}

\begin{remark}[what the analytic chain uses]\label{geoa:rem-ties}
The declaration \texttt{ChainTied.core\_of} of Lemma~\ref{geoa:tied-core} has no
callers in the archive; the main chain goes through \texttt{tiedCore\_of\_chainTied}.
On that path the conclusions use that some record satisfies the real-number clauses,
together with three parts of \texttt{TieArith}: the descriptions of $TT$ and $T_2$ as
finite sets and the identification of $\ell_\psi$ with $L'/L(1,\chi)$; the identifications of $d(u)$, $\mu_N$,
$\mu_R$ and $h(u)$ with ranks, slopes and heights of the lattices of a package are not
used after the tie is established. The lattice-theoretic content of this chapter (the
slope inequality, the frame theorem, the period lattice data) matters only for proving
that a family is tied, which happens elsewhere in the archive.
\end{remark}
